\section{Analytic consequences for finite enclosures}\label{sec:enclosures}

The preceding statements are exact. A numerical acquisition must also
control the uncertainty of every term that it uses. This section records
the analytic bounds behind such an acquisition; it does not prescribe
an implementation or assume that the boundary expression is cheaper.

\subsection{Primitive products and period products}
Suppose the cut polygon is decomposed into finitely many oriented arcs
with real parameters $s$. On each arc write $\omega=B(s)\dd s$ and let
$F$ denote the coherent primitive in Proposition~\ref{prop:stokes}.
Suppose approximations satisfy
\[
 |F-\widehat F|\le\delta_F(s),\qquad
 |B-\widehat B|\le\delta_B(s).
\]
An independent quadrature enclosure for
$(\ii/2)\sum\int\widehat F\overline{\widehat B}\dd s$ has error
$E_{\rm quad}$. Then its real part $\widehat E$ satisfies
\begin{equation}\label{eq:producterror}
 |\E-\widehat E|\le E_{\rm quad}+
 \frac12\sum\int\left(\delta_F|\widehat B|
       +|\widehat F|\delta_B+\delta_F\delta_B\right)\dd s.
\end{equation}
Indeed expand $(\widehat F+\Delta_F)
\overline{(\widehat B+\Delta_B)}$ and apply the triangle inequality.
The primitive error includes the basepoint and accumulated integral
errors; it cannot be reset independently on different boundary arcs.
Paired-edge cancellation is usable only with the corresponding common
primitive and jump enclosure.

Likewise, if $|q-\widehat q|\le\epsilon_q$ and
$|p-\widehat p|\le\epsilon_p$, then
\begin{equation}\label{eq:perioderror}
 \left|\im(\bar q p)-\im(\overline{\widehat q}\widehat p)\right|
 \le|\widehat q|\epsilon_p+|\widehat p|\epsilon_q
                         +\epsilon_q\epsilon_p.
\end{equation}
This follows by the same expansion. In a collision chart $|p|$ grows
logarithmically, so a fixed absolute uncertainty in $q$ is amplified.
A bounded exact remainder in Theorem~\ref{thm:payment} alone does not
imply a bounded acquisition error.

\subsection{Simple-root neighborhoods}
If $H(z)=(z-\alpha)U(z)$ on $|z-\alpha|\le\eta$, with
$0<\nu_1\le|U|\le M_1$, then direct polar integration gives
\begin{equation}\label{eq:cap}
 \frac{4\pi\eta}{M_1}
 \le\E(H;D(\alpha,\eta))
 \le\frac{4\pi\eta}{\nu_1}.
\end{equation}
Thus the neighborhood of a simple zero has an explicit finite payment.
Equivalently the coordinate \eqref{eq:branchregular} removes the local
singularity. An exterior floor on $H$ is never used as an interior floor
across a root. When two roots are present, the full coordinate formula
\eqref{eq:directU} is a direct integration route using exactly the same
unit and root information as the period formulas.

\subsection{A complete analytic approximation and exact deflation}
Let $H$ be holomorphic on $|z-C|\le\mathcal R$ with $|H|\le M$ there,
and let $0<R<\mathcal R$, $\rho=R/\mathcal R$. For the degree-$N$
Taylor polynomial $P_N$ about $C$, Cauchy's coefficient estimate gives,
when $N\ge2$,
\begin{align}
 \epsilon_0&:=\sup_{|z-C|\le R}|H-P_N|
 \le M\frac{\rho^{N+1}}{1-\rho},\label{eq:tay0}\\
 \epsilon_1&:=\sup_{|z-C|\le R}|H'-P_N'|
 \le\frac M{\mathcal R}\rho^N
       \left(\frac{N+1}{1-\rho}+\frac\rho{(1-\rho)^2}\right),
 \label{eq:tay1}\\
 \epsilon_2&:=\sup_{|z-C|\le R}|H''-P_N''|
 \le\frac M{\mathcal R^2}\rho^{N-1}
 \left(\frac{N(N+1)}{1-\rho}
 +\frac{(2N+1)\rho}{(1-\rho)^2}
 +\frac{\rho(1+\rho)}{(1-\rho)^3}\right).
 \label{eq:tay2}
\end{align}
For example the last formula is the sum of
$(N+1+k)(N+k)\rho^k$, using
$\sum\rho^k=(1-\rho)^{-1}$,
$\sum k\rho^k=\rho(1-\rho)^{-2}$, and
$\sum k^2\rho^k=\rho(1+\rho)(1-\rho)^{-3}$.
The first two are the corresponding zeroth and first derivatives of the
geometric tail. For the complete theta field, $M$ can be supplied by a
complete source moment; \eqref{eq:xi} is also an exact alternative at
$t=0$.

Let $\alpha,\beta$ be the exact simple roots, enclosed independently,
and define the polynomial divided difference
$Q_N(z)=P_N[\alpha,\beta,z]$. Even though $P_N$ need not vanish at the
exact roots, the complete function does vanish at them, and its divided
difference is $V(z)=H[\alpha,\beta,z]$. The elementary simplex identity proved in
Section~\ref{sec:heat} therefore yields
\begin{equation}\label{eq:deflationerror}
 |V(z)-Q_N(z)|\le\frac{\epsilon_2}{2}
 \quad(z,\alpha,\beta\in\overline{D(C,R)}).
\end{equation}
Indeed the second derivative of $H-P_N$ is integrated over a simplex of
area $1/2$; the disk is convex. This argument does not divide a field
error by $|(z-\alpha)(z-\beta)|$ near a root. It retains the exact
removable continuation and a uniform error bound. Root enclosures must
still be propagated through $Q_N$ and the coordinate map; replacing the
roots by their midpoints would not justify \eqref{eq:deflationerror}.

All these bounds apply before selecting a mesh or an arithmetic package.
They specify the error information required by either comparison route.
In particular, taking the intersection of two proved intervals is a
shared improvement available to both routes.
